\documentclass{beamer}

\usepackage[brazil]{babel}
\usepackage[latin1]{inputenc}
\usepackage[T1]{fontenc}
\usepackage{booktabs}

\hypersetup{pdfpagemode=FullScreen}

\newcommand{\estrangeiro}[1]{\textit{#1}}

\usetheme{JuanLesPins}
\usecolortheme{beaver}

\pgfdeclareimage[width=1.3cm]{university-logo}{img/ft}
\logo{\pgfuseimage{university-logo}}

%\logo{\includegraphics[width=1.3cm]{img/ft}}

\title{Uma Avalia��o de \textit{Scanners} de Vulnerabilidade na detec��o de XPath Injection}

\author[Christian]{Christian Hartung}

\institute[FT]{
  Faculdade de Tecnologia\\
  Universidade Estadual de Campinas
}
\date{22/06/2011}
\subject{Trabalho de Gradua��o Interdisciplinar}

\begin{document}

\begin{frame}[plain] % remove cabe�alhos
  \titlepage
\end{frame}

\begin{frame}
 \frametitle{Roteiro}
 \tableofcontents
\end{frame}

\section{Introdu��o}

\begin{frame}
 \frametitle{Introdu��o}
 
 \begin{itemize}[<+->]
  \item \estrangeiro{Internet}
 	\item \estrangeiro{Coud Computing}
 	\item Vulnerabilidades
 \end{itemize}
\end{frame}

\begin{frame}
 \frametitle{Casos conhecidos}
 
 \begin{itemize}[<+->]
 	\item PSN
 	  \begin{itemize}
	    \item 77 milh�es de usu�rios afetados
	    \item 12,3 milh�es com informa��es de cobran�a
	    \item \$171 milh�es
     \end{itemize}
 	\item Twitter
 	\item Sony Pictures
 \end{itemize}
\end{frame}

% Coloca um "guia" no in�cio de cada se��o
\AtBeginSection[]{\begin{frame}\frametitle{Roteiro}\tableofcontents[currentsection]\end{frame}}

\chapter{Vulnerabilidades em Aplica��es Web} \label{chap:Vulnerabilidades}
J� passou o tempo quando um \estrangeiro{firewall} era suficiente para proteger os dados de uma aplica��o \estrangeiro{web}. Algum tempo atr�s o foco dos \estrangeiro{hackers} era encontrar uma brecha no servidor e tentar explor�-la. Essa brecha podia ser uma porta que foi deixada aberta, uma senha que n�o era segura o bastante, etc.

Hoje, al�m dessas brechas no servidor hospedando a aplica��o, a pr�pria aplica��o tem sido atacada, muitas vezes com sucesso. Afinal, a aplica��o j� consulta tem acesso aos dados que esses \estrangeiro{hackers} est�o tentando obter.

Este cap�tulo vai mostrar as formas mais comumente utilizadas para realizar ataques � aplica��es \estrangeiro{web}, bem como estat�sticas sobre ataques realizados.

\section{Estado atual da seguran�a de aplica��es \estrangeiro{web}}
Atacar uma aplica��o \estrangeiro{web} pode ser t�o prejudicial quanto atacar o servidor diretamente. Se uma vulnerabilidade for encontrada e explorada, todos os dados da empresa podem estar comprometidos. A figura \ref{fig:Explorar-Vuln}, extra�da da edi��o 2010 do OWASP (\estrangeiro{Open Web Application Security Project}) Top 10\cite{OWASP:Top-10-2010}, mostra como funciona o ataque (caminho vermelho) a uma aplica��o \estrangeiro{web}. Os \emph{Thread Agents} encontraram diversas vulnerabilidades (\estrangeiro{Weakness}) na aplica��o. Essas vulnerabilidades foram exploradas e foi encontrado um caminho para dentro da empresa.

\figura{Explorar-Vuln}{Resultado da explora��o de uma vulnerabilidade}{12cm}

Um estudo anual realizado pelo WASC (\estrangeiro{The Web Application Security Consortium})\cite{WASC:Statistics} mostra que 13\% das aplica��es \estrangeiro{web} avaliadas possuem vulnerabilidades que podem ser exploradas de forma totalmente autom�tica. 49\% dessas aplica��es possuem vulnerabilidades de alto risco.

A probabilidade de encontrar uma dessas vulnerabilidades utilizando apenas \estrangeiro{scanners} � de 49\%. Se uma an�lise mais aprofundada for realizada, essa probabilidade pode subir para 96\%. O gr�fico da figura \ref{fig:Graph-Prob-Vuln2} mostra a probabilidade de encontrar vulnerabilidades utilizando v�rias formas de detec��o. Pode-se observar que � mais provavel encontrar uma vulnerabilidade que afeta o servidor.

\figura{Graph-Prob-Vuln2}{Probabilidade de uma vulnerabilidade ser detectada}{10cm}

Na figura \ref{fig:Graph-Prob-Vuln2} pode-se observar que, utilizando os dados do estudo da WASC, a probabilidade de encontrar uma vulnerabilidade no lado cliente da aplica��o utilizando \estrangeiro{scanners} de vulnerabilidade � de 63\%, enquanto que no lado servidor essa probabilidade � de 85\%. Realizando testes manuais, sem acesso ao c�digo fonte da aplica��o, essa probabilidade sobe para 71\% no lado cliente, e cai para 77\% do lado servidor. Se os testes forem realizados com acesso ao c�digo fonte, e essas probabilidades sobem para 78\% no cliente e 97\% no servidor.

\section{Principais vulnerabilidades}
Segundo o \estrangeiro{OWASP Top 10}'\cite{OWASP:Top-10-2010}, as duas vulnerabilidades mais perigosas em aplica��es \estrangeiro{web} s�o inje��o de c�digo e \estrangeiro{cross-site scripting}.

\subsection{Inje��o de c�digo}
Inje��o de c�digo (\estrangeiro{Code Injection}) ocorre quando dados sem valida��o afetam o funcionamento normal da aplica��o. 

Para explorar esse tipo de vulnerabilidade, dois objetivos devem ser atingidos\cite{SecuriFly}:
\begin{description*}
	\item[Injetar dados maliciosos na aplica��o] Isso pode ser feito utilizando um formul�rio para entrada dos dados (em uma aplica��o \estrangeiro{web}), colocando os dados maliciosos na URL (um servi�o \estrangeiro{RESTful}), manipula��o dos cabe�alhos da requisi��o, etc.
	\item[Manipular a aplica��o com esses dados]
\end{description*}

V�rias formas de manipula��o da aplica��o s�o conhecidas, e elas possibilitam obter dados protegidos, elevar o n�vel de permiss�o do usu�rio, executar comandos no sistema operacional do servidor, etc.

\subsection{Cross-Site Scripting}
Falha que permite que c�digo de outro site seja executado dentro da aplica��o. Isso pode permitir que os dados inseridos na aplica��o sejam interceptados ou conte�do indevido seja inserido na aplica��o, como propagandas.

Esta falha pode ocorrer quando dados que o usu�rio digitou s�o utilizados para desenhar a p�gina. Se o usu�rio inserir alguma \estrangeiro{tag} HTML ou algum script, ele ser� executado na p�gina\cite{SecuriFly}.

Um exemplo dessa vulnerabilidade foi o que ocorreu com o Twitter em 2010\cite{info:Twitter}.


\section{Estudo}

\begin{frame}
 \frametitle{Objetivo}
 
 \begin{itemize}[<+->]
 	\item Avaliar um conjunto de \estrangeiro{scanners} de vulnerabilidade, verificando a efic�cia na detec��o de vulnerabilidades do tipo \estrangeiro{XPath Injection}
 \end{itemize}
\end{frame}

\chapter{\estrangeiro{Scanners} de vulnerabilidades} \label{chap:scanners}
\estrangeiro{Scanners} de vulnerabilidade s�o uma forma de analisar uma aplica��o de forma autom�tica em busca de vulnerabilidades\cite{nuno}. Eles s�o utilizados como uma forma de detectar problemas gerados pela m� codifica��o da aplica��o. A maioria dos \estrangeiro{scanners} de vulnerabilidade s�o ferramentas comerciais, mas existem alguns gratuitos e at� mesmo com o c�digo fonte dispon�vel.

Segundo \cite{Tania:2010:Scanners}, essas ferramentas funcionam em tr�s fases:
\begin{description*}
	\item[Configura��o:] Nesta fase o \estrangeiro{scanner} deve ser informado sobre a aplica��o que ele varrer�, as vulnerabilidades que ser�o buscadas nessa varredura e informa��es de conex�o com o servi�o (proxy - servidor intermedi�rio entre o cliente e a aplica��o). � a �nica fase em que � necess�ria interven��o do usu�rio.
	\item[Rastreamento:] O \estrangeiro{scanner} analisa o conte�do de cada p�gina da aplica��o em busca de \estrangeiro{links} e monta um mapa de como a aplica��o est� organizada.
	\item[Explora��o:] O \estrangeiro{scanner} faz a busca pelas vulnerabilidades, fazendo simula��es da intera��o com o usu�rio (simulando um clique em um \estrangeiro{link}, o preenchimento de um formul�rio, etc.). Ao final deste processo � gerado o relat�rio com todas as vulnerabilidades detectadas.
\end{description*}

Os algor�timos utilizados para detec��o das vulnerabilidades � pr�prio de cada \estrangeiro{scanner}, o que faz com que os resultados obtidos sejam variem de \estrangeiro{scanner} para \estrangeiro{scanner}. Isso nos leva a questionar qual dos resultados � o mais correto.
\subsection{Testes realizados}

\begin{frame}
 \frametitle{Utiliza��o dos \estrangeiro{scanners}}
 
 \begin{itemize}[<+->]
 	\item Configura��o dos \estrangeiro{scanners}
 	\item An�lise da aplica��o
 	\item An�lise do relat�rio
 	\item Divertido
 \end{itemize}
\end{frame}

\begin{frame}
 \frametitle{Valida��o dos resultados}
 
 \begin{itemize}[<+->]
 	\item Verifica��o de cada resultado
 	\item Teste de Caixa-branca
 	\item Teste de Caixa-preta
 \end{itemize}
\end{frame}
\subsection{Resultados}

\begin{frame}
 \frametitle{Utiliza��o dos \estrangeiro{scanners}}
 
 \begin{table}[h]
	 \centering
	 \begin{tabular}{lrrr}
	 \toprule
	  \textbf{\estrangeiro{Scanner}} & \textbf{Encontradas} & \textbf{Falta de Cobertura} \\
	  \midrule
	  Arachni & 3 & 29 \\
	  \midrule
	  w3af & 3 & 29 \\
	  \midrule
	  WebCruiser & 5 & 27 \\
	  \bottomrule
	 \end{tabular}
	\end{table}
\end{frame}

% Remove o guia no �ltimo slide
\AtBeginSection[]{}

\chapter{Conclus�o}
Este estudo apresentou os resultados de tr�s scanners de vulnerabilidade na detec��o de uma amostra com mais de trinta vulnerabilidades comprovadas do tipo XPath Injection.

Apesar de os scanners terem conseguido n�meros melhores em outros tipos de vulnerabilidades, pode-ser perceber que n�o � t�o simples detectar XPath Injection em uma aplica��o.

Os resultados mostram que escolher um scanner de vulnerabilidade n�o � f�cil. Diversos fatores devem ser levados em considera��o, como o tempo necess�rio para scannear a aplica��o, n�mero de resultados encontrados, vulnerabilidades suportadas, etc.

Estudos futuros devem utilizar outros tipos de vulnerabilidades al�m do XPath Injection, e procurar outros scanners.

\end{document}
